\section{Discussion}\label{sec:discussion}

\subsection{What the learned rates contribute, and what they do not}
The hybrid neural ODE is the only step of the workflow that sees the data without a mechanistic hypothesis, and its output is used in three ways: to start every library mechanism from a data-informed point, to supply the targets of the symbolic regression, and to let the modeller inspect the shape of the kinetics before choosing among mechanisms (Figs.~\ref{fig:c1mu} and \ref{fig:c2virtual}). The two constraints built into the rate parameterisation, non-negativity and the vanishing of a rate without its substrate, are those shared by every candidate, so the hybrid model is mechanism-agnostic without being physically unconstrained; the ablation of Table~\ref{tab:ablation} shows that the second constraint is necessary, not cosmetic, in substrate-limited fed-batch operation. The learned rates are, however, only as accurate as the data allow. Along a single run whose states move in concert they can deviate from the true rate by tens of per cent (Section~\ref{sec:c1:hidden}), and in V15 they smoothed out the peak of the conversion rate. The workflow is therefore designed so that no decision rests on the learned rates alone: every candidate, whether from the library or from the regression, is refined against the measurements and ranked by a likelihood-based criterion, and when the best candidate fails the adequacy test the screening of Step~5 bypasses the learned rates altogether. In this sense the hybrid model is an accelerator and a generator of hypotheses, not an arbiter.

\subsection{Ambiguity is a property of the data}
The most consistent result of both case studies is that the main source of misidentification is not the fitting but the information content of the experiment. When biomass, substrate and product change monotonically and in concert, laws that attribute the slowing of a rate to different state variables make nearly the same predictions: Haldane and Contois in run~a of case study~1, product and substrate inhibition of the conversion in the BC15-type virtual run. The BIC weights and the shortlist report this ambiguity honestly rather than hiding it behind a single selected mechanism, and the remedy is experimental rather than computational. Adding two runs that place high biomass at low substrate and low biomass at high substrate raised the selection accuracy of case study~1 from \ConeAccSingle\,\% to \ConeAccJoint\,\%, and was also what made the hidden law detectable. The workflow handles several runs at no extra modelling effort---one network and one set of kinetic parameters, with one initial state per run---so that a library shortlist from one run can be used directly to plan the next one, for example with the discrimination designs of \citet{BuzziFerraris1983} or \citet{Franceschini2008}.

\subsection{Proposing mechanisms without trusting them}
Equation discovery is attractive for kinetic modelling only if it does not generate spurious mechanisms. Three design choices keep the symbolic regression conservative. First, the admissibility rules restrict the search to rate laws that a bioprocess engineer would write: non-negative, saturating denominators and rates that vanish without substrate. Second, the regression only proposes; its laws are refined against the measurements and must beat the library by a large margin on the same criterion. Third, a proposed law is accepted only if it is also adequate, so that a law that is merely the best of a poor set is not carried forward. With these rules no law was accepted in any of the \ConeLibCasesJoint{} joint and \ConeLibCasesSingle{} single-run analyses of library mechanisms, nor in the \VirtSRn{} virtual analyses, whereas the hidden law was found when the data were informative enough to require it. The regression still proved useful when nothing was accepted: it recovered the generating structure of most library laws, including the non-rational Tessier and Aiba laws through the exponential dictionary terms, and in case study~2 it consistently suggested exponential instead of hyperbolic product inhibition, an equally valid way to describe the slowing of the conversion that a modeller might wish to add to the library. Adding an accepted law to the library is also safe: it took little weight from the original mechanisms when they were the generating ones and changed none of their selections, and it remained available for later runs in which the evidence for it is weaker.

\subsection{From synthetic to real data}
Three differences between the synthetic and the real data shaped the second case study. The first is that the measurement errors are unknown. Profiling them out of the likelihood gives a ranking criterion that needs no error model, but it removes the adequacy test, and with it the safeguard against accepting a law that is merely the best of a poor set. Proposed laws on real data were therefore treated as candidates for further experiments, even where they crossed the BIC threshold. Replicate samples, or error models from the validation of the analytical methods, would restore the test. The profiled criterion also weights each species by its own fitted error, so a species that is fitted almost exactly, such as a substrate measured as zero after its depletion, carries much weight; a floor on the error standard deviations would limit this effect.

The second difference is model structure. The lumped glucose equivalent cannot represent the diauxic batch phase, in which the yeast first grows on glucose and then on the ethanol it produced, and every candidate compensated with a higher initial biomass. Because all candidates share this growth model, the error affects the ranking of the conversion mechanisms less than it affects the growth parameters, which should not be interpreted physiologically. The symbolic regression hinted at the problem, proposing growth laws that saturate more weakly than the library allows.

The third difference is run-to-run variability. Joint fitting resolved the ambiguity of single runs in case study~1 because all runs shared their kinetics by construction. The real runs agreed on the mechanism, but not on its parameters, and a joint fit with shared parameters described some runs much worse than their individual fits (Section~\ref{sec:c2:joint}). The power of joint fitting to discriminate between mechanisms therefore depends on the reproducibility of the cultures, and the shared-kinetics assumption should be tested before it is used, as was done here at no extra cost. Mixed-effects formulations, with run-specific parameters distributed around population values, are a natural way to combine runs that share a mechanism but not exact parameter values. Despite these differences, the analyses of the real runs converged on a consistent picture: every run, analysed alone, selected the mechanism with substrate and product inhibition of the conversion and cell death, and the laws proposed by the symbolic regression for every analysis contained death and described the slowing of the conversion through xylitol. Experiments that vary xylitol independently of time and xylose, for example by adding xylitol to the batch medium, would separate the proposed exponential product inhibition from the hyperbolic one of the library.


\subsection{Limitations and extensions}
Several limitations should be noted. The dictionaries of the symbolic regression are small and chosen by hand; they cover the rational and exponential laws that are customary in bioprocess kinetics, but not power laws with non-integer exponents (the Moser law of case study~1) or thresholds, and their size grows combinatorially with the number of state variables. Sequentially thresholded or mixed-integer formulations \citep{Brunton2016} would allow larger dictionaries. The workflow assumes that the mass balances are known and that the rates depend on the measured concentrations only; intracellular states, lag phases and metabolic shifts are outside its scope, as are the lumping errors of the glucose equivalent of case study~2. The adequacy test requires known measurement errors, which replicate measurements would provide for real data. The hybrid model has a few settings---the constant of the substrate factor, the integration step and the training schedule---that were fixed per case study rather than tuned. Finally, the real-data analysis rests on four runs of one strain, and the discovered laws are hypotheses to be tested in dedicated experiments rather than established mechanisms. Natural extensions are the coupling of the shortlist with model-based design of discriminating experiments, Bayesian treatment of the parameter and structural uncertainty, and the use of the accepted laws in process optimisation.
